\sectionheading{5.17 মডুলার প্রোগ্রামিং (Modular Programming)}

বড় একটি Program একসঙ্গে লিখলে তা বোঝা, পরীক্ষা করা এবং পরিবর্তন
করা কঠিন হতে পারে।

তাই বড় Program-কে কাজ অনুযায়ী কয়েকটি
\textbf{ছোট ও স্বতন্ত্র অংশে} ভাগ করে তৈরি করা হয়।
এসব অংশকে \textbf{Module} বলা হয় এবং এই পদ্ধতিই
\textbf{Modular Programming}।

%================================================
% Definition
%================================================
\booksubheading{Modular Programming কাকে বলে?}

\begin{definitionbox}

\textbf{\color{darkTeal}সংজ্ঞা}

যে Programming পদ্ধতিতে একটি বড় Program-কে নির্দিষ্ট কাজভিত্তিক
কয়েকটি \textbf{ছোট, স্বতন্ত্র ও পরিচালনাযোগ্য Module-এ ভাগ করে}
Program তৈরি করা হয়, তাকে \textbf{Modular Programming} বলে।

\end{definitionbox}

\textbf{সহজভাবে মনে রাখবে—}

\begin{center}

\textbf{
বড় Program
\(\longrightarrow\)
ছোট ছোট Module
\(\longrightarrow\)
সম্পূর্ণ Program
}

\end{center}

%================================================
% Module
%================================================
\booksubheading{Module কী?}

\begin{definitionbox}

Program-এর এমন একটি স্বতন্ত্র অংশ, যা সাধারণত
\textbf{নির্দিষ্ট একটি কাজ সম্পন্ন করে}, তাকে
\textbf{Module} বলে।

\end{definitionbox}

একটি Student Result Program-কে উদাহরণ হিসেবে এভাবে ভাগ করা যায়—

\begin{center}
\begin{adjustbox}{max width=0.98\linewidth}
\fcolorbox{darkBlue}{softBlue}{
    \parbox[c][9mm][c]{0.20\linewidth}{
        \centering\textbf{Input\\Module}
    }
}
\quad
\(\longrightarrow\)
\quad
\fcolorbox{mainTeal}{softGreen}{
    \parbox[c][9mm][c]{0.24\linewidth}{
        \centering\textbf{Calculation\\Module}
    }
}
\quad
\(\longrightarrow\)
\quad
\fcolorbox{darkBlue}{softBlue}{
    \parbox[c][9mm][c]{0.20\linewidth}{
        \centering\textbf{Result\\Module}
    }
}
\end{adjustbox}
\end{center}

এখানে—

\begin{itemize}

\item \textbf{Input Module} → নম্বর গ্রহণ করবে।

\item \textbf{Calculation Module} → Total বা Average হিসাব করবে।

\item \textbf{Result Module} → Result প্রদর্শন করবে।

\end{itemize}

সব Module একত্রে কাজ করে সম্পূর্ণ Program তৈরি করে।

%================================================
% Structure
%================================================
\booksubheading{Modular Programming-এর গঠন}

একটি বড় Problem-কে প্রথমে কয়েকটি ছোট কাজ বা Module-এ ভাগ করা হয়।

\begin{center}

\textbf{
Main Program
\(\longrightarrow\)
Module--1 + Module--2 + Module--3 + \(\cdots\)
}

\end{center}

প্রয়োজনে কোনো Module-কে আবার আরও ছোট
\textbf{Submodule}-এ ভাগ করা যেতে পারে।

\begin{examplebox}{উদাহরণ}

একটি Result Processing System-কে এভাবে ভাগ করা যেতে পারে—

\begin{itemize}

\item Student Data Module
\item Marks Calculation Module
\item Grade Determination Module
\item Result Display Module

\end{itemize}

\end{examplebox}

%================================================
% Need
%================================================
\booksubheading{Modular Programming কেন প্রয়োজন?}

বড় Program-এর সব Code একই জায়গায় লিখলে—

\begin{itemize}

\item Program জটিল হয়ে যেতে পারে।
\item Error খুঁজে পাওয়া কঠিন হতে পারে।
\item Modification করা কঠিন হয়।
\item একই Code পুনরায় ব্যবহার করা অসুবিধাজনক হতে পারে।

\end{itemize}

কাজগুলো Module-এ ভাগ করলে Program তুলনামূলকভাবে
সুশৃঙ্খল ও পরিচালনাযোগ্য হয়।

%================================================
% Complete Example
%================================================
\begin{examplebox}{একটি সহজ উদাহরণ}

ধরা যাক একটি Program-এর কাজ হলো একজন শিক্ষার্থীর তিন বিষয়ের নম্বর
নিয়ে Average এবং Result দেখানো।

Program-টিকে এভাবে ভাগ করা যেতে পারে—

\begin{tabularx}{\linewidth}{|>{\bfseries}p{0.32\linewidth}|Y|}
\hline

\rowcolor{softBlue}
\multicolumn{1}{|c|}{\textbf{Module}}
&
\multicolumn{1}{c|}{\textbf{কাজ}}
\\
\hline

Input Module
&
তিনটি নম্বর গ্রহণ
\\
\hline

Calculation Module
&
মোট ও Average নির্ণয়
\\
\hline

Result Module
&
Result নির্ধারণ
\\
\hline

Output Module
&
ফলাফল প্রদর্শন
\\
\hline

\end{tabularx}

\end{examplebox}

%================================================
% Features
%================================================
\booksubheading{Modular Programming-এর প্রধান বৈশিষ্ট্য}

\begin{tabularx}{\linewidth}{|>{\bfseries}p{0.27\linewidth}|Y|}
\hline

\rowcolor{softBlue}
\multicolumn{1}{|c|}{\textbf{বৈশিষ্ট্য}}
&
\multicolumn{1}{c|}{\textbf{ব্যাখ্যা}}
\\
\hline

বিভাজন
&
বড় Program ছোট Module-এ বিভক্ত থাকে।
\\
\hline

নির্দিষ্ট কাজ
&
প্রতিটি Module সাধারণত নির্দিষ্ট একটি কাজ সম্পন্ন করে।
\\
\hline

স্বতন্ত্রতা
&
একটি Module তুলনামূলকভাবে স্বাধীনভাবে তৈরি বা পরীক্ষা করা যায়।
\\
\hline

পুনঃব্যবহার
&
প্রয়োজন হলে কোনো Module পুনরায় ব্যবহার করা যেতে পারে।
\\
\hline

সহজ Testing
&
Module আলাদাভাবে পরীক্ষা করা সুবিধাজনক।
\\
\hline

সহজ Maintenance
&
নির্দিষ্ট অংশ পরিবর্তন বা সংশোধন করা সহজ।
\\
\hline

\end{tabularx}

%================================================
% Advantages
%================================================
\booksubheading{Modular Programming-এর সুবিধা}

\begin{compacttwocolpoints}

\item বড় ও জটিল Program সহজে পরিচালনা করা যায়।

\item Program বোঝা সহজ হয়।

\item Error শনাক্ত ও Debugging করা সহজ হয়।

\item নির্দিষ্ট Module পরিবর্তন করলে পুরো Program পরিবর্তনের
প্রয়োজন কমে।

\item একই Module অন্য জায়গায় পুনরায় ব্যবহার করা যেতে পারে।

\item একাধিক Programmer বিভিন্ন Module নিয়ে আলাদাভাবে কাজ করতে পারে।

\item Testing ও Maintenance সহজ হয়।

\end{compacttwocolpoints}

%================================================
% Limitations
%================================================
\booksubheading{Modular Programming-এর সীমাবদ্ধতা}

\begin{itemize}

\item Program-কে সঠিকভাবে Module-এ ভাগ করার জন্য আগে পরিকল্পনা প্রয়োজন।

\item Moduleগুলোর মধ্যে সঠিক যোগাযোগ বা Data আদান-প্রদান নিশ্চিত করতে হয়।

\item খুব ছোট Program-কে অপ্রয়োজনীয়ভাবে বেশি Module-এ ভাগ করলে
Program-এর Structure জটিল হতে পারে।

\end{itemize}

\begin{rememberbox}

\textbf{\color{deepOrange}মনে রাখবে}

Modular Programming-এর উদ্দেশ্য শুধু বেশি Module তৈরি করা নয়;
বরং \textbf{কাজ অনুযায়ী যৌক্তিকভাবে Program ভাগ করা}।

\end{rememberbox}

%================================================
% Module and Function
%================================================
\booksubheading{Module ও Function-এর সম্পর্ক}

একটি নির্দিষ্ট কাজকে আলাদা করে বাস্তবায়নের জন্য Programming-এ
\textbf{Function} ব্যবহার করা যায়।

তাই Function Modular Programming বাস্তবায়নের একটি গুরুত্বপূর্ণ
উপায় হতে পারে।

উদাহরণ—

\begin{center}

\textenglish{\ttfamily input\_data()}
\(\longrightarrow\)
Input গ্রহণ

\vspace{0.5mm}

\textenglish{\ttfamily calculate()}
\(\longrightarrow\)
হিসাব করা

\vspace{0.5mm}

\textenglish{\ttfamily show\_result()}
\(\longrightarrow\)
Result দেখানো

\end{center}

Python Function পরবর্তী অংশে বিস্তারিত আলোচনা করা হবে।

%================================================
% Comparison
%================================================
\booksubheading{Modular ও Non-Modular Program-এর তুলনা}

\begin{tabularx}{\linewidth}{|>{\bfseries}p{0.22\linewidth}|Y|Y|}
\hline

\rowcolor{softBlue}
\multicolumn{1}{|c|}{\textbf{বিষয়}}
&
\multicolumn{1}{c|}{\textbf{Modular Program}}
&
\multicolumn{1}{c|}{\textbf{Non-Modular/একক Program}}
\\
\hline

গঠন
&
ছোট ছোট Module-এ বিভক্ত
&
অধিকাংশ কাজ একই Program অংশে থাকে
\\
\hline

বোঝা
&
তুলনামূলক সহজ
&
বড় হলে কঠিন হতে পারে
\\
\hline

Testing
&
Module আলাদাভাবে পরীক্ষা করা যায়
&
পুরো Program একসঙ্গে পরীক্ষা করতে হতে পারে
\\
\hline

Debugging
&
Error-এর অংশ শনাক্ত করা সহজ
&
তুলনামূলক কঠিন
\\
\hline

Modification
&
নির্দিষ্ট Module পরিবর্তন করা যায়
&
পরিবর্তনে অন্য অংশ প্রভাবিত হতে পারে
\\
\hline

Reusability
&
বেশি
&
তুলনামূলক কম
\\
\hline

Maintenance
&
সহজ
&
বড় Program-এ তুলনামূলক কঠিন
\\
\hline

\end{tabularx}

%================================================
% Extra Concept
%================================================
\booksubheading{Main Module ও Submodule — সংক্ষিপ্ত ধারণা}

বড় Program-এ যে অংশ অন্য Module-গুলোর কাজ সমন্বয় করে,
তাকে সাধারণভাবে \textbf{Main Module} হিসেবে ভাবা যায়।

কোনো Module-কে আবার আরও ছোট অংশে ভাগ করলে সেগুলোকে
\textbf{Submodule} বলা যায়।

\begin{center}

\textbf{
Main Module
\(\longrightarrow\)
Module
\(\longrightarrow\)
Submodule
}

\end{center}

এটি মূল ধারণার তুলনায় কম গুরুত্বপূর্ণ;
Program-এর Structure বোঝার জন্য সংক্ষেপে জানলেই যথেষ্ট।

%================================================
% Final Reminder
%================================================
\begin{rememberbox}

\textbf{\color{deepOrange}মনে রাখবে}

\begin{itemize}

\item \textbf{Modular Programming} =
বড় Program-কে কাজভিত্তিক ছোট Module-এ ভাগ করা।

\item \textbf{Module} সাধারণত একটি নির্দিষ্ট কাজ সম্পন্ন করে।

\item একটি Module আবার \textbf{Submodule}-এ বিভক্ত হতে পারে।

\item Modular Programming-এ
\textbf{Testing, Debugging, Modification ও Maintenance} সহজ হয়।

\item Module-এর \textbf{Reusability} একটি গুরুত্বপূর্ণ সুবিধা।

\item Function ব্যবহার করে Modular Programming বাস্তবায়ন করা যায়।

\item খুব বেশি অপ্রয়োজনীয় Module তৈরি করাও ভালো Design নয়।

\item দ্রুত মনে রাখার কৌশল—

\begin{center}

\textbf{
Divide
\(\longrightarrow\)
Develop
\(\longrightarrow\)
Test
\(\longrightarrow\)
Combine
}

\end{center}

\end{itemize}

\end{rememberbox}

\small

%================================================
% 5.18 Program Design Techniques
%================================================
